\documentclass[11pt,a4paper]{article}
\usepackage[margin=25mm]{geometry}
\usepackage[T1]{fontenc}
\usepackage{lmodern}
\usepackage{amsmath,amssymb}
\usepackage{enumitem}
\usepackage{graphicx}
\usepackage{booktabs}
\usepackage[hidelinks]{hyperref}
\setlist[enumerate]{leftmargin=*,itemsep=0.6em,parsep=0.2em}
\setlist[itemize]{itemsep=0.2em}
\setlength{\emergencystretch}{3em}
\newcommand{\R}{\mathbb{R}}

\title{Multi-Robot Exploratory Coverage Path Planning in Unknown Indoor Dynamic Environments with MARL\\[0.3em]
\large Problem Formulation}
\author{}
\date{}

\begin{document}
\maketitle

\begin{figure}[t]
    \centering
    \includegraphics[width=0.6\linewidth]{images/robot_coverage_vector.pdf}
    \caption{2D floor of the MR-ECPP task. Robots, equipped with 2D LiDAR
    sensors, are drawn as coloured discs; pedestrians as grey ellipses.}
    \label{fig:floor}
\end{figure}

\noindent We consider a cooperative \textbf{Multi-Robot Exploratory Coverage
Path Planning (MR-ECPP)} problem. A team of robots must physically pass over
every reachable part of an indoor floor that none of them knows in advance,
while avoiding walls, each other and, optionally, pedestrians. The robots
therefore have to \emph{explore} the floor and \emph{cover} it at the same
time.

The approach rests on four ideas:
\begin{itemize}
 \item \textbf{Local knowledge only.} Each robot builds its own map of what
       it has seen and covered, and merges it with nearby teammates when they
       are in communication range, an approach inspired by
       \cite{xiao2020a_distributed}. No robot ever has the true map.
 \item \textbf{A learned policy.} Each robot is driven by a neural network
       that outputs velocity commands directly. It is trained with
       Multi-Agent Reinforcement Learning (MAPPO) under centralized training
       with decentralized execution (CTDE).
 \item \textbf{A good starting point.} Before reinforcement learning, the
       policy can be initialized by Imitation Learning (IL) from a classical
       Boustrophedon Cellular Decomposition (BCD) expert.
 \item \textbf{A safety net.} A classical A*{}+DWA controller takes over when
       the learned behaviour stops making progress, inspired by
       \cite{carvalho2025deep}.
\end{itemize}
The problem is formulated by the following items.

%==========================================================================
\section*{Environment and robots}

\begin{enumerate}[label=\textbf{\arabic*.}]

\item \textbf{Workspace and layouts.}
The floor is a $12\times 8$~m rectangle divided into rooms, corridors and
doorways (Fig.~\ref{fig:floor}). Interior walls are aligned with the grid and
are at least one cell thick. At every episode the floor plan is drawn from a
bank of $64$ procedurally generated layouts, so the policy learns a general
strategy instead of memorizing one floor. The size of the floor is fixed, but
its layout is unknown to the robots.

\item \textbf{Continuous motion, discrete coverage.}
Robots move in continuous space, while maps and coverage use a square grid
of cell size $0.5$~m ($24\times16$ cells). When a robot moves, the cell under
its centre becomes \emph{covered}, together with each neighbouring cell that
the robot body enters by at least $80\%$ of its radius. This threshold avoids
counting cells that the robot only grazes. Coverage is permanent within an
episode. The episode ends when all reachable free cells are covered, or after
$T_{\max}=4000$ steps ($400$~s).

\item \textbf{Robot team.}
The team is $\mathcal I=\{1,\dots,N\}$ with $N=3$. Robots are homogeneous:
they have the same sensors, motion limits and capabilities. They start at
random collision-free positions.

\item \textbf{Kinematics.}
Each robot is a disc of radius $r=0.2$~m with pose $q_i=(x_i,y_i,\theta_i)$
and differential-drive (unicycle) kinematics
\[
 \dot x_i=v_i\cos\theta_i,\qquad
 \dot y_i=v_i\sin\theta_i,\qquad
 \dot\theta_i=\omega_i,
\]
integrated with a time step $\Delta t=0.1$~s, which is also the decision
period of the policy.

\item \textbf{Continuous actions.}
At each step the policy outputs $a_i\in(-1,1)^2$ (a tanh-squashed Gaussian),
mapped to
\[
 v_i\in[0,v_{\max}],\quad \omega_i\in[-\omega_{\max},\omega_{\max}],
 \qquad v_{\max}=1~\text{m/s},\ \omega_{\max}=1~\text{rad/s}.
\]
The robot drives forward only and steers by turning. MARL approaches to
MR-ECPP with continuous actions are currently a gap in the literature.

\item \textbf{Collisions.}
Collisions are \emph{soft}. A move that would hit a wall, a robot or a
pedestrian is rejected: the robot stays where it is for that step and
receives a penalty, but it keeps operating and the episode continues. This
lets the policy learn from its mistakes without ending the episode.

\item \textbf{Pedestrians (optional).}
Pedestrians walk along straight segments in random directions and stop
instead of stepping towards a nearby robot. During training a
\emph{ghost robot} curriculum is used: at first pedestrians ignore the robots,
so the robots must learn to avoid them on their own; as the policy improves,
pedestrians increasingly react. At test time they always react. The main
experiments use no pedestrians.

\item \textbf{Sensing and localization.}
Each robot knows its initial pose in a common global frame and has
\emph{perfect odometry}, so it always knows where it is, without drift.
Each robot carries a 2D LiDAR with $L=70$ rays over $360^\circ$ and
$r_L=5$~m range. The LiDAR detects walls, other robots and pedestrians, and
is used both to discover the floor and to avoid obstacles.

\end{enumerate}

%==========================================================================
\section*{Map memory and communication}

\begin{enumerate}[label=\textbf{\arabic*.},resume]

\item \textbf{Individual map belief.}
\label{item:individual_map_belief}
Each robot $i$ keeps its own map $M_i$, a grid of the same size as the
floor that starts empty and is filled during the episode. The map tells the
robot what it knows: where the walls are, which areas it has already
covered, and which areas still need to be covered. This is what lets the
robot decide where to go without access to the true map.

Every cell of $M_i$ is in one of four states:
\[
 \textit{unknown},\quad
 \textit{occupied},\quad
 \textit{free--uncovered},\quad
 \textit{free--covered}.
\]
The map is updated in three ways:
\begin{itemize}
 \item \textit{Sensing}: cells seen by the LiDAR become known, either as
       free--uncovered or as occupied (walls);
 \item \textit{Covering}: a free cell becomes covered only when the robot
       body passes over it, so seeing a cell is not the same as covering it;
 \item \textit{Sharing}: robots in communication range exchange their maps,
       so each one also learns what its teammates have seen and covered.
\end{itemize}
Information is only ever added: once a cell is known it stays known, and
once it is covered it stays covered.

\item \textbf{Communication.}
Two robots communicate when their distance is below
$R_{\mathrm{comm}}=3$~m. Connected robots share their positions and merge
their maps cell by cell, keeping the most informative state of each cell,
so after an exchange they share the same belief. Exchange happens at every
step and only between direct neighbours; information still spreads through
the team over successive steps, from robot to robot. When disconnected,
robots keep sensing, mapping and acting on their own. They may then cover
the same area twice, which is the price of not relying on communication.

\end{enumerate}

%==========================================================================
\section{Actor observation}

\begin{enumerate}[label=\textbf{\arabic*.},resume]

\item \textbf{Observation frame.}
At each step robot $i$ builds its observation only from its own sensing and
its own map. One frame holds:
\begin{itemize}
 \item \textbf{Pose}: $(x/W,\ y/H,\ \cos\theta,\ \sin\theta)$, where $W$ and
       $H$ are the floor dimensions.
 \item \textbf{Last executed command}: $(v/v_{\max},\ \omega/\omega_{\max})$.
 \item \textbf{Robot id}: $i\in\{1,\dots,N\}$.
 \item \textbf{Teammates in communication range}: the $2$ nearest teammates
       within $R_{\mathrm{comm}}$, each as $(\mathrm{id}_j,\ \Delta x_j,\
       \Delta y_j)$, their relative position in metres. An empty slot is all
       zeros, and $\mathrm{id}=0$ marks it as empty.
 \item \textbf{History}: a validity flag and the relative position of the
       centre of the previously visited cell, followed by the current sweep
       direction $d\in\{-1,0,1\}^2$, the direction of the robot's last step
       between two adjacent new cells. They tell the robot where it comes
       from and which way it is sweeping.
 \item \textbf{LiDAR}: $L=70$ ranges, normalized by $r_L=5$~m.
 \item \textbf{Local crop} (extending \cite{hasan2025efficient}): a
       $3\times7\times7$ window of the robot's map
       (item~\ref{item:individual_map_belief}) around the robot. It works
       like a small camera looking down on the robot. Its inner $5\times5$
       core is centred on the robot's cell and has three binary channels:
       \begin{itemize}
        \item \textit{obstacle}: a wall cell the robot knows about;
        \item \textit{covered}: a cell covered according to the robot's map;
        \item \textit{robots}: a cell where a teammate in communication range
              stands.
       \end{itemize}
       Cells beyond the floor boundary are marked as obstacles. The one-cell
       ring around the core summarizes the rest of the floor: each side of
       the ring holds, for each channel, the average over all the cells
       beyond that side, i.e.\ the fraction of known walls, of covered cells
       and of cells with a teammate. A side with no cell inside the floor
       takes the beyond-the-floor values (obstacle $1$, covered $0$, robots
       $0$), and the corners are $0$. The core gives detail where the robot
       is; the ring tells it, coarsely, in which direction work is left.
\end{itemize}

\begin{figure}[t]
    \centering
    \makebox[\textwidth][c]{%
        \includegraphics[width=1.1\textwidth]{images/NN_architecture.pdf}%
    }
    \caption{Actor policy architecture.}
    \label{fig:nn_architecture}
\end{figure}

A frame therefore has $18$ scalar entries, $70$ LiDAR ranges and a
$3\times7\times7$ crop. The full map is \emph{not} an actor input: it stays
inside the robot and reaches the policy only through the crop. No target
cell, waypoint or route is given: the actor must decide where to go from the
crop, its history and the ring.

\item \textbf{Short-term memory.}
A single frame does not show how the robot is moving. Two variants give the
actor a short-term memory:
\begin{itemize}
 \item \textbf{Observation stacking} (main variant): the last $N_s=5$
       frames, $0.5$~s of history, are given to the actor in time order.
 \item \textbf{GRU memory} (alternative): the actor reads one frame per step
       and a recurrent GRU unit (Fig.~\ref{fig:nn_architecture}) carries the
       history in its internal state.
\end{itemize}

\item \textbf{Actor network.}
The LiDAR scan is read by a circular 1D convolutional network, so the two
ends of the scan stay adjacent, and the crop by a small 2D convolutional
network. Their features are joined with the scalar entries and passed to
fully connected layers that output the mean of the action distribution. Its
spread does not depend on the observation and is bounded, so exploration
never collapses or explodes.

\item \textbf{Parameter sharing.}
All robots use the same actor, with the same weights. The robots are
identical and the floors are unknown, so no robot should specialize in a
particular role or area. Sharing also means that the experience of every
robot trains the same policy, so the policy learns $N$ times faster than with
separate networks.

\end{enumerate}

%==========================================================================
\section{Global state for the critic}

\begin{enumerate}[label=\textbf{\arabic*.},resume]

\item \textbf{A privileged view.}
The critic is used only during training, so it may read information that no
robot has: the true floor plan, the true team coverage and the exact state of
the whole team. A single critic network $V_\phi$, shared by all robots, is
evaluated once per robot on an \emph{agent-specific} global state
$s_t^i = f_i(s_t)$: the same global state $s_t$, seen from the point of view
of robot $i$ \cite{yu2022surprising}. Own cell, own motion and own task
context are marked separately from the rest of the team, so the critic gives
a different value to each robot, consistent with the per-robot rewards.

\item \textbf{Structure of the state.}
The input of the critic for robot $i$ is a short history of \emph{global
frames}:
\[
  s_t^i = \big(g_{t-K+1}^i,\ \dots,\ g_t^i\big),
  \qquad
  g_t^i = \big(\mathbf{M}_t^i,\ \mathbf{z}_t^i\big),
  \qquad K=5.
\]
Each frame has two parts: a stack of maps $\mathbf{M}_t^i$ for everything
that has a position on the floor, and a vector $\mathbf{z}_t^i$ for
everything that does not. Stacking frames lets the critic see trends, such
as a robot that keeps going back and forth, much as the actor does with its
own frames. At the start of an episode the missing frames are left empty.

\item \textbf{Global maps $\mathbf{M}_t^i$.}
Grids over the whole floor that tell the critic how much work is left, where
it is, and who is where:
\begin{itemize}
 \item \textbf{Walls}: the true floor plan, including unexplored parts, so
       the critic knows how large the task really is.
 \item \textbf{Coverage}: the cells the team has actually covered. Together
       with the walls, it shows the remaining work.
 \item \textbf{Self}: the cell of robot $i$, which places it relative to the
       remaining work.
 \item \textbf{Teammates}: the cells of all other robots, wherever they are.
       Unlike the actor, the critic is not limited by communication range.
 \item \textbf{Visit counts}: one map per robot, recording how often it has
       entered each cell. They distinguish efficient coverage from robots
       going over the same cells again, and show which robot is doing it.
\end{itemize}

\item \textbf{Global vector $\mathbf{z}_t^i$.}
Information that is not spatial or is hard to read from a grid:
\begin{itemize}
 \item \textbf{Own motion}: pose and velocity of robot $i$, placed first so
       the critic knows which robot it is evaluating.
 \item \textbf{Team motion}: pose and velocity of all robots, and positions of
       the pedestrians.
 \item \textbf{Task context of robot $i$}: the internal quantities that
       shape its rewards and decide whether the recovery controller takes
       over: elapsed episode time, revisit streak and steps without progress,
       recovery state and goal, current sweep and lane state, and the current
       weight of the revisit penalty. None of these appears in the maps, yet
       each one changes the future return.
 \item \textbf{Previously visited cell} of every robot, which adds a sense of
       direction to the team's motion.
\end{itemize}

\item \textbf{Critic network.}
A small convolutional network reads the stacked maps and summarizes the
spatial situation. The summary is joined with the global vector, and fully
connected layers turn both into the value $V_\phi(s_t^i)$: an estimate of
the return robot $i$ can still expect from time $t$ on
(Fig.~\ref{fig:critic_network}). The critic is only a training signal: once
training is over it is discarded, and each robot acts using only its own
observation.

\begin{figure}[h]
    \centering
    \includegraphics[width=0.9\linewidth]{images/critic_network.pdf}
    \caption{Centralized critic architecture.}
    \label{fig:critic_network}
\end{figure}

\end{enumerate}

%==========================================================================
\section{Learning: MAPPO with CTDE}

\begin{enumerate}[label=\textbf{\arabic*.},resume]

\item \textbf{Multi-agent problem.}
The robots learn coverage by trial and error with multi-agent reinforcement
learning (MARL). The task is a cooperative partially observable Markov game
$\langle \mathcal{N}, \mathcal{S}, \{\mathcal{A}^i\}, P, \{r^i\},
\{\Omega^i\}, \gamma \rangle$: a Dec-POMDP in which each robot receives its
own reward $r^i_t$. Two difficulties arise: \emph{partial observability},
because each robot sees only its own observation $o^i_t$, and
\emph{non-stationarity}, because all robots learn at the same time, so from
the point of view of one robot the environment keeps changing.

\item \textbf{CTDE.}
At execution each robot runs its own copy of the \emph{actor}
$\pi_\theta(a^i_t\mid o^i_t)$ and decides only from its own observation,
with no central computer. During training in simulation, the \emph{critic}
$V_\phi$ reads the privileged global state and judges the robots' actions.
Since it sees the whole team, changes in the teammates' behaviour appear in
its input instead of looking like noise.

\item \textbf{PPO.}
Comparing the critic's value with what actually happened gives the
advantage $\hat{A}^i_t$, computed with GAE: positive if an action turned out
better than expected, negative if worse. The actor is updated with
PPO~\cite{schulman2017proximal}:
\[
  L(\theta) = \mathbb{E}\Big[\min\big(\rho^i_t\,\hat{A}^i_t,\
  \operatorname{clip}(\rho^i_t,\,1-\epsilon,\,1+\epsilon)\,\hat{A}^i_t\big)\Big],
  \qquad
  \rho^i_t = \frac{\pi_\theta(a^i_t\mid o^i_t)}{\pi_{\theta_{\text{old}}}(a^i_t\mid o^i_t)} .
\]
Good actions become more likely and bad ones less likely, while clipping the
ratio $\rho^i_t$ limits how far a single update can move the policy. An
entropy bonus keeps the robots exploring.

\item \textbf{MAPPO.}
MAPPO~\cite{yu2022surprising} brings PPO to CTDE. All robots share one actor
and one centralized critic, which is evaluated once per robot. Each robot's
advantage is computed only from its own rewards and values. This keeps
credit assignment clear: the robot that covered new cells is credited, and
an idle teammate is not.

\item \textbf{Training.}
Training alternates between collecting experience, with $32$ parallel
environments for $128$ steps each on randomly drawn layouts, and updating
actor and critic on it. Four mechanisms keep learning stable:
\begin{itemize}
 \item the actor stops updating when the policy changes too much, measured
       by the KL divergence between old and new policy;
 \item the critic uses a Huber loss, which is precise on small errors and
       robust to large ones, so rare large rewards do not disrupt it;
 \item rewards are scaled and observations normalized;
 \item steps driven by the recovery controller are excluded from PPO, since
       the actor did not choose them, and are imitated instead with a
       separate supervised loss.
\end{itemize}
The hyperparameters are listed in Table~\ref{tab:ppo_hparams}.

\begin{table}[h]
\centering
\begin{tabular}{ll}
\toprule
$\gamma$ / GAE $\lambda$ & $0.99$ / $0.95$ \\
Clipping $\epsilon$ / entropy coefficient & $0.2$ / $0.01$ \\
Learning rate actor / critic & $10^{-4}$ / $10^{-3}$, linear decay \\
Epochs / minibatches & $4$ / $8$ \\
KL threshold / reward scale & $0.02$ / $0.05$ \\
\bottomrule
\end{tabular}
\caption{PPO hyperparameters.}
\label{tab:ppo_hparams}
\end{table}

\end{enumerate}

%==========================================================================
\section{Reward function}

\begin{enumerate}[label=\textbf{\arabic*.},resume]

\item \textbf{Structure.}
Each robot $i$ receives its own reward $r_t^i$ at every step. It is the sum
of terms, each encouraging or discouraging one specific behaviour:
\[
  r_t^i =
  \underbrace{r^{\mathrm{cov}} + r^{\mathrm{done}}}_{\text{task}}
  + \underbrace{r^{\mathrm{seq}} + r^{\mathrm{lane}} + r^{\mathrm{axis}}}_{\text{sweep pattern}}
  + \underbrace{r^{\mathrm{prog}} + r^{\mathrm{loit}} + r^{\mathrm{rev}}}_{\text{efficiency}}
  + \underbrace{r^{\mathrm{spr}}}_{\text{coordination}}
  + \underbrace{r^{\mathrm{col}} + r^{\mathrm{fb}} - \tau}_{\text{safety, time}} .
\]
The reward is computed by the simulator and used only during training, so it
may use the true coverage. Terms that say where a robot \emph{should} go
(progress and lane shaping) use only the robot's own map, so the actor is
never asked to react to information it cannot see. Below, $C_t\in[0,1]$ is
the fraction of the floor covered by the team, $T_{\max}$ the episode length,
and $\mathbb{1}[P]$ the indicator function, equal to $1$ if condition $P$
holds and $0$ otherwise. The robot index $i$ is dropped when clear from
context. Parameter values are listed in Table~\ref{tab:reward_params}.

%--------------------------------------------------------------------------
\item \textbf{Task terms.}
\begin{itemize}
 \item \textbf{New coverage.}
 \[
   r^{\mathrm{cov}}_t = \alpha\,(1 + g\,C_{t-1})\; n_t ,
 \]
 where $n_t$ is the number of cells newly covered by the robot at this step.
 When two robots reach the same new cell at the same time, only one of them
 is credited. This is the only term that directly measures the objective.
 The factor $(1+g\,C_{t-1})$ makes late cells worth up to three times more
 than early ones: late cells are usually the hardest to reach, such as
 isolated corners or cells left behind, and without this factor the
 incentive to finish would fade exactly when it is needed most.

 \item \textbf{Completion.}
 \[
   r^{\mathrm{done}}_t = B\,\Big(1 + \beta\,\big(1 - \tfrac{t}{T_{\max}}\big)\Big)\,
   \mathbb{1}[\text{floor fully covered at } t],
 \]
 given to \emph{every} robot when the last cell is covered. The shared bonus
 makes the team succeed or fail together, so a robot has a reason to help
 finish an area it did not start. The time factor makes an early finish worth
 up to twice as much as a late one, rewarding speed and not only success.
\end{itemize}

%--------------------------------------------------------------------------
\item \textbf{Sweep-pattern terms.}
Coverage alone does not care \emph{how} the floor is covered. These terms
favour the Boustrophedon pattern of a lawn mower: long straight lanes joined
by one-cell shifts with a reversal of direction, which covers a region
without gaps and with few re-entries.

Both terms look at the sequence of \emph{new} cells a robot covers. Let
$c_{\text{last}}$ be the last new cell the robot covered and $\delta_t$ the
step from $c_{\text{last}}$ to the new cell covered now. The new cell is
\emph{adjacent} if $\|\delta_t\|_1 = 1$, i.e.\ it touches $c_{\text{last}}$
on one of its four sides. The step $\delta_t$ is compared with two reference
directions:
\begin{itemize}
  \item $d$, the previous step: $\delta_t = d$ means the robot keeps going
        straight;
  \item $\ell$, the return direction: after a straight run of at least two
        cells along a lane, a one-cell sideways shift sets $\ell$ to the
        opposite of that lane's direction. $\delta_t = \ell$ means the robot
        goes back along the new lane, completing the U-turn. When there was
        no shift, $\ell$ is not defined and this case never occurs.
\end{itemize}

\begin{itemize}
 \item \textbf{Sequential coverage: how the robot moves.}
 \[
   r^{\mathrm{seq}}_t = \mathbb{1}[\|\delta_t\|_1 = 1]\,
   \big(b_{\mathrm{seq}} + b_{\mathrm{str}}\,\mathbb{1}[\delta_t = d]
         + b_{\mathrm{bou}}\,\mathbb{1}[\delta_t = \ell]\big).
 \]
 An adjacent new cell earns $b_{\mathrm{seq}}$; going straight earns
 $b_{\mathrm{str}}$ more; the U-turn on the new lane earns $b_{\mathrm{bou}}$
 more. A new cell that is not adjacent earns nothing. This term rewards
 covering in connected strips instead of jumping from place to place, but it
 only looks at the shape of the move: it cannot tell whether turning was a
 good idea.

 \item \textbf{Lane shaping: when the robot should turn.}
 Let $p$ be the direction the robot is expected to follow: $\ell$ right
 after a lane shift, $d$ otherwise. The cell \emph{ahead} is
 $c_{\text{last}} + p$. It is \emph{open} if, according to the robot's own
 map, it is neither a wall nor already covered. Then
 \begin{align*}
    r^{\mathrm{lane}}_t ={} & b_{\mathrm{ls}}\,\mathbb{1}[\text{new cell straight ahead}] \\
    & {}+ b_{\mathrm{lt}}\,\mathbb{1}[\text{new cell sideways},\ \text{ahead not open}] \\
    & {}- c_{\mathrm{lb}}\,\mathbb{1}[\text{robot leaves the lane},\ \text{ahead open}].
 \end{align*}
 In words: keep going straight while the lane still has work; turn sideways
 only when the lane is finished or blocked; leaving a lane that is still open
 is penalized, because the robot would have to come back later for the
 leftover piece.
\end{itemize}
The two terms complement each other. If a robot turns sideways onto a new
cell while the cell ahead is still open, $r^{\mathrm{seq}}_t$ rewards the
move (the cell is adjacent) but $r^{\mathrm{lane}}_t$ penalizes it (the lane
was not finished). If the cell ahead were a wall, both would reward the turn.
Sequential coverage teaches \emph{how} to move; lane shaping teaches
\emph{when} to change lane.

\begin{itemize}
 \item \textbf{Axis alignment.}
 \[
   r^{\mathrm{axis}}_t = -\,c_{\mathrm{ax}}
   \Big(\frac{e_t}{45^\circ}\Big)^{2}
   \min\!\Big(1,\ \frac{\|\Delta \mathbf{x}_t\|}{v_{\max}\Delta t}\Big),
 \]
 where $e_t\in[0^\circ,45^\circ]$ is the angle between the robot's
 displacement $\Delta\mathbf{x}_t$ and the nearest grid axis. Lanes follow
 rows and columns of cells, and diagonal motion crosses cells only partially,
 leaving slivers uncovered. Squaring the error barely penalizes small
 drifts, and the speed factor means that turning on the spot costs nothing.
\end{itemize}

%--------------------------------------------------------------------------
\item \textbf{Efficiency terms.}
\begin{itemize}
 \item \textbf{Progress towards work.}
 Let $D^i(\mathbf{x})$ be the distance, in cells, from position $\mathbf{x}$
 to the nearest cell that robot $i$ \emph{believes} is free and uncovered,
 measured along routes through cells it knows to be free. When no such cell
 is reachable, the nearest exploration frontier (a known free cell next to
 unknown space) is used instead. Then
 \[
   r^{\mathrm{prog}}_t = \lambda_p\,\big(D^i(\mathbf{x}_{t-1}) - D^i(\mathbf{x}_t)\big).
 \]
 Once the robot's surroundings are finished, the coverage reward is zero and
 gives no direction. This term fills that gap by pulling the robot towards
 the nearest remaining work. It follows the idea of potential-based shaping
 \cite{ng1999policy}: as long as the map does not change, it sums to zero
 over any closed loop, so it cannot be farmed by driving back and forth.

 \item \textbf{Loitering.}
 \[
   r^{\mathrm{loit}}_t = -\,c_{l}\;\mathbb{1}[n_t = 0]\;
   \Big(1 - \mathrm{clip}\big(\tfrac{\Delta D_t}{\Delta D_{\max}},0,1\big)\Big),
 \]
 where $\Delta D_t$ is the progress of the step and $\Delta D_{\max}$ the
 progress of a full-speed step straight towards work. A step that neither
 covers a new cell nor approaches work is wasted, so parking, spinning or
 wandering over covered cells is charged. Driving at full speed towards work
 costs nothing.

 \item \textbf{Re-coverage.}
 \begin{align*}
    r^{\mathrm{rev}}_t &= -\,w(\phi_t)\,\min(k_t,\,k_{\max})\;\mathbb{1}[\text{entered a covered cell}], \\[0.5em]
    w(\phi) &= c_{\mathrm{rv}}\big[\eta + (1-\eta)(1-\phi)^2\big],
    \qquad
    \phi_t = \max\!\big(C_t,\ t/T_{\max}\big),
 \end{align*}
 where $k_t$ counts the consecutive entries into covered cells since the
 robot last covered a new one. Crossing covered cells is the main source of
 inefficiency, and the growing streak punishes long detours more than a
 single unavoidable crossing. The weight $w(\phi)$ is high early on, when
 fresh work is nearby, and decays towards $\eta\,c_{\mathrm{rv}}$ as coverage
 or time advances, because near the end the remaining cells are scattered
 and reaching them \emph{requires} crossing covered ones.
\end{itemize}

%--------------------------------------------------------------------------
\item \textbf{Coordination term.}
\[
  r^{\mathrm{spr}}_t =
  -\,w_s \!\!\sum_{j\in\mathcal{N}_i^{\mathrm{comm}}}\!\! \Big[1 - \frac{d_{ij}}{r_s}\Big]_+
  \;-\; w_{\mathrm{los}}(t) \!\!\sum_{j\in\mathcal{N}_i^{\mathrm{los}}}\!\! \Big[1 - \frac{d_{ij}}{r_{\mathrm{los}}}\Big]_+ ,
  \qquad
  w_{\mathrm{los}}(t) = w_{\mathrm{los}}\Big[1-(1-\eta_{\mathrm{los}})\min\!\big(1,\tfrac{t}{T_d}\big)\Big],
\]
where $d_{ij}$ is the distance between robots $i$ and $j$,
$[\cdot]_+=\max(0,\cdot)$, $\mathcal{N}_i^{\mathrm{comm}}$ the teammates in
the robot's communication slots and $\mathcal{N}_i^{\mathrm{los}}$ the
teammates in direct line of sight. Robots close to each other compete for
the same cells, so the team covers fewer cells per step than it could. The
two repulsions push the team to spread out: the first acts on teammates
known through communication, the second on any teammate the robot can see,
and is strongest at the start of the episode, when the robots must leave the
same area. No room assignment or negotiation is involved: coordination
emerges from these terms and from the shared completion bonus.

%--------------------------------------------------------------------------
\item \textbf{Safety and time terms.}
\begin{itemize}
 \item \textbf{Collisions}:
 $r^{\mathrm{col}}_t = -\kappa_w\,\mathbb{1}[\text{wall}] - \kappa_r\,\mathbb{1}[\text{robot}]
 - \kappa_h\,\mathbb{1}[\text{pedestrian}]$. The penalty grows with the
 severity of the contact, so hitting a person is the worst case.
 \item \textbf{Recovery activation}:
 $r^{\mathrm{fb}}_t = -c_{\mathrm{fb}}\,\mathbb{1}[\text{recovery controller activated}]$,
 charged once, to the action that triggered it. Without this cost the
 policy could rely on the recovery controller instead of learning to get out
 of trouble by itself.
 \item \textbf{Time}: a constant $-\tau$ per step, a small, steady pressure to
 finish quickly.
\end{itemize}

\begin{table}[h]
\centering
\caption{Reward parameters.}
\label{tab:reward_params}
\begin{tabular}{llll}
\toprule
\textbf{Symbol} & \textbf{Value} & \textbf{Symbol} & \textbf{Value} \\
\midrule
$\alpha$, $g$                        & $10$, $2$          & $c_l$                                 & $0.2$ \\
$B$, $\beta$                         & $200$, $1$         & $c_{\mathrm{rv}}$, $\eta$, $k_{\max}$ & $6$, $0.025$, $3$ \\
$b_{\mathrm{seq}}$, $b_{\mathrm{str}}$, $b_{\mathrm{bou}}$ & $6$, $4$, $8$ & $w_s$, $r_s$       & $2$, $3$\,m \\
$b_{\mathrm{ls}}$, $b_{\mathrm{lt}}$, $c_{\mathrm{lb}}$    & $3$, $3$, $5$ & $w_{\mathrm{los}}$, $r_{\mathrm{los}}$ & $1.5$, $5$\,m \\
$c_{\mathrm{ax}}$                    & $0.5$              & $\eta_{\mathrm{los}}$, $T_d$          & $0.3$, $1500$ steps \\
$\lambda_p$                          & $2$                & $\kappa_w$, $\kappa_r$, $\kappa_h$    & $2$, $5$, $10$ \\
$\tau$                               & $0.02$             & $c_{\mathrm{fb}}$                     & $10$ \\
\bottomrule
\end{tabular}
\end{table}

\end{enumerate}

%==========================================================================
\section{Imitation learning initialization}

\begin{enumerate}[label=\textbf{\arabic*.},resume]

\item \textbf{Why imitation first.}
A systematic sweep is hard to discover from rewards alone: at the start the
policy moves randomly and rarely produces a clean lane. The actor is
therefore first trained to imitate a classical expert, and reinforcement
learning then refines this starting behaviour.

\item \textbf{BCD expert.}
For each layout the expert splits the free space into rectangles, sweeps
each rectangle in Boustrophedon lanes along its longer side, and chains the
rectangles into a single tour of all free cells. The tour is cut into $N$
pieces of similar length, one per robot. While driving, each robot heads for
the first cell of its piece that its own map does not hold as covered.

\item \textbf{Behaviour cloning with DAgger.}
Rollouts are collected while the robots execute either the expert's command
or the actor's own, with the share of expert commands decreasing over time.
The expert labels every visited state, so the actor also learns to recover
from situations it reaches by itself. The actor regresses the expert's
velocity commands, with more weight on decisions that change the previous
command, such as turns, lane shifts and stops. A label is used only when the
expert's target lies inside the $5\times5$ core of the actor's crop: choosing
a distant target requires planning over the whole map, which is left to
reinforcement learning and to the recovery controller.

\item \textbf{Transition to reinforcement learning.}
Before MAPPO starts, the critic is warmed up on the imitation policy, with the
actor frozen, so that PPO starts from a critic that already knows the value
of the policy it is updating. MAPPO then fine-tunes the policy from this
checkpoint, optionally keeping a decaying imitation term at the start.

\end{enumerate}

%==========================================================================
\section{Recovery controller (fallback)}

\begin{enumerate}[label=\textbf{\arabic*.},resume]

\item \textbf{Trigger.}
Following the hybrid idea of \cite{carvalho2025deep}, a robot switches to a
classical controller when the learned policy is clearly stuck: after $7$
consecutive entries into covered cells, or after $70$ steps ($7$~s) without
covering a new cell.

\item \textbf{Planning and execution.}
A* selects the nearest uncovered cell that the robot's own map shows as
reachable, routing only through known free cells. DWA follows the route,
respecting acceleration limits and the current LiDAR clearance. If DWA stops
making progress, a sequence of slow velocity commands ($\le 0.4$~m/s), each
checked against the latest LiDAR, drives the robot along the route. The
controller keeps control until the robot itself covers the goal cell; if a
teammate covers it first, a new goal is planned. It never uses the true map.

\item \textbf{Consistency with learning.}
Steps driven by the recovery controller are not samples of the actor, so
they are excluded from the PPO policy update; the critic still learns from
all steps. Collision-free recovery commands train the actor through a
separate imitation loss, so the policy gradually learns to get out of
trouble by itself. Results are reported both with the controller (assisted)
and without it (policy only).

\end{enumerate}

%==========================================================================
\section{Evaluation}

\begin{enumerate}[label=\textbf{\arabic*.},resume]

\item \textbf{Metrics.}
Success (full coverage) rate; robot--robot, robot--wall and
robot--pedestrian collision rates; timeout rate; completion time of
successful episodes; re-coverage ratio (cell visits / distinct covered
cells); and recovery usage. Policy-only and assisted runs are compared on
the same random seeds.

\item \textbf{To be defined.}
Ablations on communication, local crop vs.\ full map, observation stacking
vs.\ GRU, IL initialization and recovery controller; generalization to
unseen layouts and to other team sizes, map sizes and pedestrian densities.
Graph Attention Networks \cite{velivckovic2017graph}, reported to improve
coordination in multi-agent exploration \cite{chiun2025marvel}, are a
candidate for the actor--critic architecture.

\end{enumerate}

%==========================================================================
\section*{Open questions}

\begin{enumerate}
\item When two robots are very close, their observations are almost
identical: how can they learn not to execute the same action? The robot id
and the teammate slots break the symmetry only partially.
\end{enumerate}

\bibliographystyle{IEEEtran}
\bibliography{references}

\end{document}
